\textbf{Phần 1. Trắc nghiệm lựa chọn}
\setcounterpageref{ex}{0}

\begin{ex} %[3D1N04]
Kết quả của phép nhân $2 \times 7$ là
\choice
{$12$}
{$16$}
{\True $14$}
{$18$}
\loigiai{
$2 \times 7 = 14$.\\
\textbf{Đáp án của câu hỏi là C}
}
\end{ex}

\begin{ex} %[3D1N04]
Kết quả của phép nhân $5 \times 8$ là
\choice
{$35$}
{\True $40$}
{$45$}
{$30$}
\loigiai{
$5 \times 8 = 40$.\\
\textbf{Đáp án của câu hỏi là B}
}
\end{ex}

\begin{ex} %[3D1H04]
Kết quả của phép chia $18 : 2$ là
\choice
{$8$}
{$7$}
{$6$}
{\True $9$}
\loigiai{
$18 : 2 = 9$.\\
\textbf{Đáp án của câu hỏi là D}
}
\end{ex}

\begin{ex} %[3D1H04]
Kết quả của phép chia $45 : 5$ là
\choice
{\True $9$}
{$8$}
{$7$}
{$6$}
\loigiai{
$45 : 5 = 9$.\\
\textbf{Đáp án của câu hỏi là A}
}
\end{ex}

\begin{ex} %[3D1H04]
Điền số thích hợp vào chỗ chấm: $5 \times \dots = 35$.
\choice
{$6$}
{$8$}
{$5$}
{\True $7$}
\loigiai{
$35 : 5 = 7$.\\
\textbf{Đáp án của câu hỏi là D}
}
\end{ex}

\begin{ex} %[3D1H04]
Điền số thích hợp vào chỗ chấm: $\dots : 2 = 8$.
\choice
{$14$}
{$18$}
{\True $16$}
{$12$}
\loigiai{
$8 \times 2 = 16$.\\
\textbf{Đáp án của câu hỏi là C}
}
\end{ex}

\begin{ex} %[3D1V04]
Tính giá trị biểu thức: $2 \times 9 + 32$.
\choice
{$48$}
{\True $50$}
{$52$}
{$40$}
\loigiai{
$2 \times 9 + 32 = 18 + 32 = 50$.\\
\textbf{Đáp án của câu hỏi là B}
}
\end{ex}

\begin{ex} %[3D1V04]
Tính giá trị biểu thức: $50 : 5 - 6$.
\choice
{$3$}
{$5$}
{\True $4$}
{$6$}
\loigiai{
$50 : 5 - 6 = 10 - 6 = 4$.\\
\textbf{Đáp án của câu hỏi là C}
}
\end{ex}

\begin{ex} %[3D1V04]
Mỗi bàn học có 2 bạn học sinh. Hỏi 9 bàn học như thế có tất cả bao nhiêu học sinh?
\choice
{$16$ học sinh}
{\True $18$ học sinh}
{$20$ học sinh}
{$14$ học sinh}
\loigiai{
$2 \times 9 = 18$ học sinh.\\
\textbf{Đáp án của câu hỏi là B}
}
\end{ex}

\begin{ex} %[3D1V04]
Có 35 bông hoa cắm đều vào 5 lọ. Mỗi lọ có số bông hoa là
\choice
{\True $7$ bông}
{$6$ bông}
{$8$ bông}
{$9$ bông}
\loigiai{
$35 : 5 = 7$ bông hoa.\\
\textbf{Đáp án của câu hỏi là A}
}
\end{ex}

\begin{ex} %[3D1V04]
Một can đựng $5\text{ l}$ dầu. Hỏi 6 can như thế đựng bao nhiêu lít dầu?
\choice
{$25\text{ l}$}
{$35\text{ l}$}
{$20\text{ l}$}
{\True $30\text{ l}$}
\loigiai{
$5 \times 6 = 30\text{ l}$.\\
\textbf{Đáp án của câu hỏi là D}
}
\end{ex}

\begin{ex} %[3D1C04]
Có 20 cái kẹo chia đều cho 2 bạn An và Bình. Sau đó An ăn hết 3 cái kẹo. Hỏi An còn lại bao nhiêu cái kẹo?
\choice
{\True $7$ cái}
{$6$ cái}
{$8$ cái}
{$5$ cái}
\loigiai{
Mỗi bạn nhận: $20 : 2 = 10$ cái kẹo.\\
An còn lại: $10 - 3 = 7$ cái kẹo.\\
\textbf{Đáp án của câu hỏi là A}
}
\end{ex}

\textbf{Phần 2. Trắc nghiệm đúng sai}
\setcounterpageref{ex}{0}

\begin{ex} %[3D1H04]
Xét tính đúng sai của các phép nhân, chia:
\choiceTF
{\True $2 \times 8 = 16$.}
{\True $40 : 5 = 8$.}
{$5 \times 0 = 5$.}
{\True $5 \times 0 = 0$.}
\loigiai{
\begin{itemize}
    \item $2 \times 8 = 16$: Đúng.
    \item $40 : 5 = 8$: Đúng.
    \item Số nào nhân 0 cũng bằng 0 $\neq 5$. Ý 3 Sai, Ý 4 Đúng.
\end{itemize}
\begin{itemize}
    \item a) \textbf{Đúng}.
    \item b) \textbf{Đúng}.
    \item c) \textbf{Sai}.
    \item d) \textbf{Đúng}.
\end{itemize}
}
\end{ex}

\begin{ex} %[3D1V04]
Có 30 quyển vở chia đều cho 5 bạn. Xét tính đúng sai:
\choiceTF
{\True Mỗi bạn nhận được 6 quyển vở.}
{\True Cả 5 bạn nhận tất cả 30 quyển vở.}
{Nếu chia cho 2 bạn thì mỗi bạn nhận được 12 quyển vở.}
{\True Nếu chia cho 2 bạn thì mỗi bạn nhận được 15 quyển vở.}
\loigiai{
\begin{itemize}
    \item $30 : 5 = 6$ quyển: Đúng.
    \item Tổng 30 quyển: Đúng.
    \item $30 : 2 = 15$ quyển $\neq 12$. Ý 3 Sai, Ý 4 Đúng.
\end{itemize}
\begin{itemize}
    \item a) \textbf{Đúng}.
    \item b) \textbf{Đúng}.
    \item c) \textbf{Sai}.
    \item d) \textbf{Đúng}.
\end{itemize}
}
\end{ex}

\textbf{Phần 3. Trả lời ngắn}
\setcounterpageref{ex}{0}

\begin{ex} %[3D1H04]
Tính tích: $2 \times 6$.
\shortans{12}
\loigiai{
\textbf{Đáp số:} $12$\\
$2 \times 6 = 12$.
}
\end{ex}

\begin{ex} %[3D1H04]
Tính thương: $25 : 5$.
\shortans{5}
\loigiai{
\textbf{Đáp số:} $5$\\
$25 : 5 = 5$.
}
\end{ex}

\begin{ex} %[3D1V04]
Tính giá trị biểu thức: $5 \times 9 - 15$.
\shortans{30}
\loigiai{
\textbf{Đáp số:} $30$\\
$45 - 15 = 30$.
}
\end{ex}

\begin{ex} %[3D1V04]
Mỗi đôi đũa có 2 chiếc đũa. Hỏi 8 đôi đũa có bao nhiêu chiếc đũa?
\shortans{16}
\loigiai{
\textbf{Đáp số:} $16$\\
$2 \times 8 = 16$ chiếc đũa.
}
\end{ex}

\begin{ex} %[3D1V04]
Có 20 bạn học sinh xếp thành các hàng, mỗi hàng có 5 bạn. Hỏi xếp được bao nhiêu hàng?
\shortans{4}
\loigiai{
\textbf{Đáp số:} $4$\\
$20 : 5 = 4$ hàng.
}
\end{ex}

\begin{ex} %[3D1C04]
Tính giá trị biểu thức: $2 \times 4 \times 5$.
\shortans{40}
\loigiai{
\textbf{Đáp số:} $40$\\
$8 \times 5 = 40$.
}
\end{ex}
